\documentclass[10pt,a4paper]{amsart}
\usepackage[margin=19mm]{geometry}
\usepackage{amsmath,amssymb,mathtools}
\usepackage[scr=rsfs,cal=euler]{mathalfa}
\usepackage{xfrac}
\usepackage[unicode,hidelinks,bookmarks=false]{hyperref}
\newcommand{\ie}{i.e.\ }
\newcommand{\cf}{cf.\ }
\newcommand{\resp}{respectively\ }
\newcommand{\Subsection}{\subsection}
\providecommand{\nobreakdash}{\nobreak\hbox{-}}
\setlength{\emergencystretch}{2em}
\multlinegap0pt
\usepackage{amssymb}
\usepackage[shortlabels]{enumitem}
\usepackage[normalem]{ulem}
\usepackage{xcolor}
\newtheorem{lemma}{Lemma}
\newtheorem{proposition}{Proposition}
\usepackage{cleveref}
\crefname{lemma}{Lemma}{Lemmas}
\crefname{proposition}{Proposition}{Propositions}
\title{Classification and nonexistence \\ for $t$-structures on \\ derived categories of schemes}
\author{Alexander Clark, Pat Lank, Kabeer Manali Rahul, Chris J. Parker}
\date{Source: arXiv:2404.08578v5 (CC BY 4.0)}
\begin{document}
\maketitle

\noindent\emph{Source and licence.} Classification and nonexistence \\ for $t$-structures on \\ derived categories of schemes. Version arXiv:2404.08578v5 (CC BY 4.0), distributed by arXiv under the Creative Commons Attribution 4.0 International licence, http://creativecommons.org/licenses/by/4.0/, as recorded in the arXiv licence metadata for this exact version and shown on the abstract page https://arxiv.org/abs/2404.08578v5. Full licence notice: \texttt{licenses/arxiv-2404.08578-CC-BY-4.0.txt}.\par\smallskip

\noindent\textit{Editorial scope.} Counted unit: the article's proposition prop:reduction\_to\_affines\_stalks (source0.tex lines 611--617). The article's complete original proof of it is delivered with it (source0.tex lines 619--629, one \texttt{proof} environment of 11 lines). No mathematics is written, repaired or re-derived: the statement and the proof are the article's own lines, in the article's own notation, and no retained line is substituted or re-flowed.  Retained uncounted as the source-local context the proof needs: lines 592--598.  Nothing was omitted from any retained span, so the package carries no omission record.  No retained line contains a citation, so the wrapper carries no bibliography and no bibliography entry.  No retained line contains a figure environment, an image-inclusion command or drawing code, so no image is delivered and the package declares no asset dependency.  Editorial formatting only, in addition: an AMS \texttt{amsart} preamble that restates the article's own declarations and macros used by the retained lines, namely the environments \texttt{lemma}, \texttt{proposition} on separate counters, together with the standard packages those lines need.  The retained environments print in this document as Lemma 1, Proposition 1 in order of appearance, whereas the article prints the same items with the numbering of its own full document.  The complete native source of the article as distributed in the arXiv e-print for this exact version is bundled unchanged as \texttt{sources/157622/source0.tex}.
\begin{lemma}\label{lem:global_to_open_perfect_supported}
    Let $Z$ be a closed subset of $X$. Then there are essentially dense functors:
    \begin{enumerate}
        \item $j^\ast \colon \operatorname{Perf}_Z (X)\to \operatorname{Perf}_{Z\cap U} (U)$ where $j\colon U \to X$ is an open immersion
        \item $s^\ast \colon \operatorname{Perf}_Z (X) \to \operatorname{Perf}_{Z\cap \operatorname{Spec}(\mathcal{O}_{X,p})} (\mathcal{O}_{X,p})$ where $s\colon \operatorname{Spec}(\mathcal{O}_{X,p}) \to X$ is the natural morphism for $p$ in $X$.
    \end{enumerate}
\end{lemma}

\begin{proposition}\label{prop:reduction_to_affines_stalks}
   Let $Z$ be a closed subset of $X$. Suppose $\phi$ is a Thomason filtration on $X$ supported in $Z$. Assume the $t$-structure on $D_{\operatorname{qc}}(X)$ associated to $\phi$ restricts to a $t$-structure on $\operatorname{Perf}_Z (X)$. 
   \begin{enumerate}
       \item If $j \colon U \to X$ is an open immersion from an affine scheme, then the $t$-structure associated to the Thomason filtration $\phi\cap U$ on $D_{\operatorname{qc}}(U)$ restricts to $\operatorname{Perf}_{Z\cap U} (U)$.
       \item For any point $p$ in $Z$, the $t$-structure associated to the Thomason filtration $\phi_p$ on $D_{\operatorname{qc}}(\mathcal{O}_{X,p})$ restricts to $\operatorname{Perf}_{Z\cap \operatorname{Spec}(\mathcal{O}_{X,p})}(\mathcal{O}_{X,p})$.
   \end{enumerate}
\end{proposition}

\begin{proof}
   We will prove $(1)$ as $(2)$ is shown in a similar manner. Choose an object $P$ in $\operatorname{Perf}_{Z\cap U} (U)$. From \Cref{lem:global_to_open_perfect_supported}, there exists $P^\prime$ in $\operatorname{Perf}_Z (X)$ such that $P$ is a direct summand of $j^\ast P^\prime$ in $\operatorname{Perf}_{Z\cap U}( U)$. Consider the truncation triangle in $D_{\operatorname{qc}}(X)$:
   \begin{displaymath}
       A \to P^\prime \to B \to A[1]
   \end{displaymath}
   where $A$ is in $\mathcal{A}$ and $B$ is in $\mathcal{B}[-1]$. Our hypothesis tells us that $A,B$ belong to $\operatorname{Perf}_Z (X)$ as $P^\prime$ is in $\operatorname{Perf}_Z (X)$ and $\mathcal{A}$ restricts to $\operatorname{Perf}_Z(X)$. The restriction of $\phi$ to $D_{\operatorname{qc}}(U)$ has associated $t$-structure $(j^\ast \mathcal{A},j^\ast \mathcal{B})$, and we have a corresponding truncation triangle for $j^\ast P^\prime$:
   \begin{displaymath}
       j^\ast A \to j^\ast P^\prime \to j^\ast B \to j^\ast A[1].
   \end{displaymath}
   where $j^\ast A$ is in $j^\ast \mathcal{A}$ and $j^\ast B$ is in $j^\ast \mathcal{B}[-1]$, both belonging to $\operatorname{Perf}_{Z\cap U} (U)$. There exist maps $f\colon P \to j^\ast P^\prime$ and $g\colon j^\ast P^\prime \to P$ whose composition is the identity. Denote the truncation functors of the $t$-structure $(j^\ast \mathcal{A},j^\ast \mathcal{B})$ on $D_{\operatorname{qc}}(U)$ by $\tau^{\leq 0} \colon D_{\operatorname{qc}}(U) \to \mathcal{A}$ and $\tau^{\geq 1} \colon D_{\operatorname{qc}}(U) \to \mathcal{B}$. Applying the connective truncation functors $\tau^{\leq 0}, \tau^{\geq 1}$ tells us that $\tau^{\leq 0} P$ is a direct summand of $j^\ast A$ and $\tau^{\geq 1} P$ is a direct summand of $j^\ast B$. In other words, the truncations of $P$ with respect to $(j^\ast \mathcal{A}, j^\ast \mathcal{B})$ belong to $\operatorname{Perf}_{Z\cap U} (U)$.
\end{proof}

\end{document}
